\uuid{uo9i}
\titre{Matrice/Puissance/Formule du binôme}
\auteur{hanani ; mzari}
\organisation{exo7}
\licence{CC BY-NC-SA 4.0}
\datecreate{2020-01-20}
\niveau{L1}
\module{Algèbre}
\chapitre{Matrice}
\sousChapitre{Propriétés élémentaires, généralités}
\difficulte{3}
\theme{Calcul matriciel, Matrice, Puissance, Formule du binôme}
\exo7qcmid{lille-2.matrices.017}

\contenu{
\begin{question}
\qclassification{108.01}
\qtags{motcle=Matrice/Puissance/Formule du binôme}
Soit $a,b\in \Rr$, $A=\left(\begin{array}{rcc}a&1&b\\0&a&2\\ 0&0&a\\
\end{array}\right)$ et  $N= A-aI$, où $I= \left(\begin{array}{rcc}
1&0&0\\0&1&0\\ 0&0&1\\ \end{array}\right)$. Quelles sont les assertions vraies ?
\begin{answers}
\good{$N^k = 0$, pour tout entier $k\ge 3$}.
\bad{On ne peut pas appliquer la formule du binôme pour le calcul de $A^n$}.
\good{Pour tout entier $n \ge 2,$ $\displaystyle A^n =a^nI+na^{n-1}N+\frac{n(n-1)}{2}a^{n-2}N^2 $}.
\good{Pour tout entier $n \ge 2,$ $A^n=\left(\begin{array}{rcc}
a^n&na^{n-1}&na^{n-1}b+n(n-1)a^{n-2}\\0&a^n&2na^{n-1}\\ 0&0&a^n\\ \end{array}\right)$}.
\end{answers}
\vskip2mm
\begin{explanations} On a : $N=\left(\begin{array}{rcc}
0&1&b\\0&0&2\\ 0&0&0\\ \end{array}\right)$,  $N^2=\left(\begin{array}{rcc}0&0&2\\0&0&0\\ 0&0&0\\
\end{array}\right)$ et $N^k=0, $ pour tout $k\ge 3$.
\vskip0mm
On a : $A= N+aI$. Comme le produit des matrices $N$ et $aI$ est commutatif, on peut appliquer la formule du binôme pour le calcul des puissances de $A$.
\end{explanations}
\end{question}
}
